\section{Results, comparison, and logical meaning}
\label{sec:overview}

There is no single useful ordering of ``simplest'' Diophantine forms.
The number of unknowns, degree, number of equations, arithmetic circuit
size, and size of fixed coefficients are distinct resources. We therefore
give explicit forms on several parts of that tradeoff, with a common
fully specified coefficient tuple.

Throughout, $\N=\{0,1,2,\ldots\}$ and $\Zpos=\{1,2,\ldots\}$.
All additional letters introduced in formulas are polynomial
abbreviations, not existentially quantified unknowns. A circuit operation
means one binary addition, subtraction, or multiplication. Squaring and
multiplication by a fixed coefficient count; shared registers and fixed
leaves do not. Powers in displayed formulas always have fixed
nonnegative integer exponents and are expanded into multiplication gates
in the delivered circuits.

\begin{theorem}[Explicit consistency forms]
\label{thm:main}
With the four fixed native ports specified in
Section~\ref{sec:complemented-compiler}, the archive defines polynomials
$F_9$, $F_9^+$, $F_{10}$, and $F_{11}$ such that ordinary mathematics
formalizable in ZFC proves
\begin{align*}
 \Con(\mathrm{ZFC})
 &\ \longleftrightarrow\ \neg\exists\mathbf x\in\N^9\ F_9(\mathbf x)=0,\\
 &\ \longleftrightarrow\ \neg\exists\mathbf x\in\Zpos^9\ F_9^+(\mathbf x)=0,\\
 &\ \longleftrightarrow\ \neg\exists\mathbf x\in\Z^{10}\ F_{10}(\mathbf x)=0,\\
 &\ \longleftrightarrow\ \neg\exists\mathbf x\in\Zpos^{11}\ F_{11}(\mathbf x)=0.
\end{align*}
The last line also has an equivalent two-equation form in the same
eleven unknowns. The first two lines use the parameterized Jones ratio
theorems. The third additionally uses the stated Lucas and interpolation
results of Sun~\cite{sun2026}. Every line uses the inherited ordinary
compiler and literal K215 calibration at the pinned revision~\cite{repo}.
\end{theorem}

\begin{table}[htbp]
\centering
\small
\begin{tabular}{@{}llllr@{}}
\toprule
Form & Domain & Unknowns & Exact degree(s) & Operations\\
\midrule
Quadratic batching & $\N$ & 9 & $166\,177\,600$ & 868\\
Positive translation & $\Zpos$ & 9 & $166\,177\,600$ & 864\\
Sun companion & $\Z$ & 10 & $449\,607\,016$ & 773\\
Index absorption & $\Zpos$ & 11 & $1154$ & 123\\
Two-equation form & $\Zpos$ & 11 & $2,1152$ & 119\\
Shifted two-equation form & $\Zpos$ & 11 & $2,1152$ & 118\\
Translated product & $\Zpos$ & 11 & $1140$ & 129\\
Translated shifted system & $\Zpos$ & 11 & $2,1138$ & 124\\
\bottomrule
\end{tabular}
\caption{New forms in the delivered archive. The certificate circuit for
the nine-variable degree proof is intentionally a different evaluator.}
\label{tab:results}
\end{table}

\subsection{What improves over the pinned repository}

The nine-variable degree previously reached
\[
 92324\cdot5^{32}+296
 =2\,149\,585\,634\,469\,985\,961\,914\,062\,796.
\]
The new degree is smaller by a factor of approximately
$1.294\cdot10^{19}$. The 864-operation positive circuit also improves
the separate historical 1473-operation positive circuit by $41.34\%$;
the degree-optimized package~37 circuit uses 1500 operations. These are
separate baseline comparisons, not an invented baseline combining the
best features of different evaluators.

The ten-integer baseline in package~31 has degree
\[
 126496\cdot5^{31}+27064
 =589\,042\,901\,992\,797\,851\,562\,527\,064
\]
and 1442 operations. Quadratic batching reduces the degree by
approximately $1.310\cdot10^{18}$ and the circuit cost by $46.39\%$.
The domain remains unrestricted integers.

The eleven-positive construction removes one supplied unknown from the
direct twelve-positive Pell-norm family. It trades degree and arithmetic
for that reduction: the repository still has a twelve-positive degree
1022 construction with 124 operations and a separate 108-operation
degree 1081 construction. The present work does not improve the
historical 67-operation single equation or the 59-operation system.
It also leaves intact the usefulness of the 31-positive quadratic system
and its quartic sum of squares when low degree is the primary goal.

The native ports now use $339{,}476$ raw bytes rather than $506{,}664$,
a $32.998\%$ reduction. Their largest bit length falls from
$1{,}016{,}946$ to $681{,}374$. This is a reduction in native compiler
constants; it is not a claim about the size of a fully expanded
$166$-million-degree polynomial.

\subsection{The three ideas}

The first change is to preserve the system's quadratic structure.
Multiplying a quadratic coefficient kernel by the square of a digit code
can select a residual by one coefficient. A Sidon set gives unambiguous
quadratic monomials with maximum position 799. Disjoint translated
blocks select all 22 residuals simultaneously. A single binary carry
test enforces every target coefficient to be zero, so there is no need
to take a sum of squares before entering the few-variable interface.
This reduces both the degree and the work needed to evaluate the source
code. It also removes the old enormous Jones ceiling constant: the
nine-variable source can use simply $b=\epsilon+2$ and
$B=\kappa b^2$.

The second change is arithmetic rather than combinatorial. The
eleven-variable norm initially recovers only a weakened divisibility
$E\mid H(k-R-1)$. Cancelling $H$ here would be wrong. Instead we
recover the main Pell index first, place the remaining index in an
interval shorter than $E/\gcd(E,H)$, and deduce the exact equality
$2n=R+1$. This restores the full divisibility and reconstructs the
deleted positive quotient. The proof explicitly avoids a predecessor
bound that required the deleted quotient.
A further positive translation of the same carrier lowers degree 1154
to 1140 at the cost of six operations; it is supplied as a separate
optional form.

The third change uses a whole-label complemented encoding. An occupied
label $x$ is represented by $2^\ell-1-x$. A selected tile contributes
$x+1$, so one guarded binary field checks the entire label by requiring
its low $\ell$ bits to vanish. This replaces one guard field per bit
by one per label while retaining the occupied high bit needed by the
native compiler. We recompute the complete coefficient tuple and a
compatible packed-rank controller; changing only a few constants would
not suffice.

\subsection{What the consistency claim means}

Here $\Con(\mathrm{ZFC})$ means that no finite syntactic ZFC refutation
exists, using the effective proof presentation calibrated in the
repository. The fixed K215 machine satisfies the inherited theorem
\[
 \text{K215 halts on blank input}
 \quad\longleftrightarrow\quad
 \text{a finite ZFC refutation exists}.
\]
The named repository endpoint is
\texttt{KEndpoint215.blank\_halts\_iff\_zfc\_refutation}.
The compiler and arithmetic reductions connect a halt with an integer
zero, in both directions. Negation gives Theorem~\ref{thm:main}.

This is a theorem about equivalence. It does not prove either side of
the consistency statement. In particular it creates no conflict with
the second incompleteness theorem. All witnesses are ordinary finite
integers; none is a variable exponent, an unspecified enumeration index,
or an oracle for consistency.

The archive separates ordinary proofs, formal source endpoints, and
finite checks. We have not rebuilt the entire upstream proof-assistant
development or formalized the new batching and absorption arguments in
its kernel. The checks establish the finite facts claimed in their
receipts; they do not establish an unbounded equivalence by testing.
There is also no claim that nine is minimal for this particular
consistency predicate, or that the Sidon set or operation counts are
globally optimal.
